\SourceLeaf{23}
weil die Summe von allen $=0$) sein, die nicht $\equiv 0$ ist. Ebenso wenig
aber können die Grössen $\alpha, \beta, \gamma$ alle unter einander
incongruent sein, weil sie sonst, gleichviel in welcher Ordnung, den Zahlen
$0, 1, 2$ congruent sein und also ihre Summe durch $3$ theilbar sein
müsste; es wird folglich unter den Differenzen $\alpha-\beta$,
$\beta-\gamma$, $\gamma-\alpha$ gewiss Eine, aber auch nur Eine sein, die
$\equiv 0$ ist. Es gibt mithin drei Fälle, die einander ausschliessen:

\[
\text{I.} \quad \text{Wenn } \beta \equiv \gamma \quad(\mathrm{mod.}\ 3)
\]
\[
\text{II.} \quad \text{Wenn } \gamma \equiv \alpha \quad(\mathrm{mod.}\ 3)
\]
\[
\text{III.} \quad \text{Wenn } \alpha \equiv \beta \quad(\mathrm{mod.}\ 3).
\]

Welcher dieser drei Fälle statthat, hängt von $x$ ab; wir werden daher nach
Maassgabe dieser drei Fälle sämmtliche durch $p$ nicht theilbare Zahlen in
drei Classen theilen, in die erste nämlich diejenigen setzen, wo der
erste; in die zweite die, wo der zweite; in die dritte die, wo der dritte
Fall stattfindet.

\textls[100]{Erstes Theorem:} Alle nach dem Modulus $p$ congruente Zahlen gehören in
Einerlei Classe.

\textls[100]{Zweites Theorem:} Gleiche, aber mit entgegengesetzten Zeichen afficirte
Zahlen gehören in Eine Classe.

Hieraus ist klar, dass man bloss die Zahlen $1, 2, 3, \ldots,
\frac{1}{2}(p-1)$ zu classificiren braucht, um sofort alle übrigen
classificiren zu können.

\textls[100]{Drittes Theorem:} Die Zahlen $(3n-m)x$, $2mx$ gehören in zwei auf einander
folgende Classen.

\textls[100]{Viertes Theorem:} Also wenn $\dfrac{2n-m}{3m} \equiv f$, also

\[
1+f+ff \equiv 0 \quad(\mathrm{mod.}\ p), \qquad 1 \equiv f^3 \quad(\mathrm{mod.}\ p),
\]
\noindent so gehören $x$ und die Reste von $fx$, $ffx$ in drei in umgekehrter
Ordnung auf einander folgende Classen.

